% Pursuing Stacks (cotes 134-2 to 134-8): archive cote 134-3, site page 87 (PDF page 88), typescript page 87.
% typescript. Draft compared: ps2.tex:4548-4619.
% AI-transcribed by Claude Opus 5.5 (claude-opus-5-5, adaptive thinking, effort high), Message Batches API, batch ps1 (msgbatch_01Rv5hCaNTtUt4JHx1yqoHPL), 2026-11-12; not yet checked by eye.
\scanpage{134-3/87}
\mspage{\hand{87}}

a morphism of intervals\hand{,} namely compatible with endpoints. (However in prac\unsure{-}%
% NOTE: right edge of the scan cuts the end of this line after "prac"
tical terms (12) is the more ready-to-use criterion,  because in most
cases I will be in A and therefore  $i_{\hand{!}}(I) = i(I)$,  and (12) will be more
or less trivial,  for instance because I will be chosen so that $i(I)$ is
either $\hand{\Delta}_1 = \hand{\{0\}\to\{1\}}$  or it's barycentric subdivision \hand{$\begin{tikzcd}[row sep=-10pt,column sep=small] \{0\} \ar[dr] & \\ & \{0,1\} \\ \{1\} \ar[ur] & \end{tikzcd}$}.)  The
same argument shows that (iii) implies that A is a test category\hand{,} namely
the Lawvere element $L_A$ in $A\hat{}$ is a homotopy interval in $A\hat{}$\hand{,} namely is asphe-%
ric over $e_{A\hat{}}$ .  This being so,  we \xout{deducexthatxi$_A$xinducesxanx} only got to
prove (under the assumption (iii))  that \inserted{$i_Ai^*$}(11) induces an autoequivalence
of (Hot),  and more specifically a functor isomorphic to the \unsure{identify|identity}
functor\hand{.} \inserted{This we achieved} by comparing directly $A_{/C}$ to C  by the functor  $(a,p)\hand{\mapsto} p(e_a)$
(13) \qquad $A_{/C} \hand{\longrightarrow} C$ ,
where  $e_a$ is the final element in $i(a)$.  It is enough to check this is a
weak equivalence for any C,  a fortiori that it is aspheric,  and the usu\xout{\unsure{a}}%
al asphericity criterion reduces us to showing it is a weak equivalence
when C has a final object,  namely that in this case  $A_{/C} = A_{/i^*(C)}$  is
aspheric,  i.e. $i^*(\unsure{C})$ is aspheric in $A\hat{}$.  This was achieved by an extreme-%
ly simple homotopy argument,  using a homotopy  in (Cat)
\[ \hand{\Delta}_1 \times C \hand{\to} C \]
between the identity functor in C and the constant functor with value $e_C$ ;
(cf p. 62).

I feel like reserving the appe\hand{\unsure{ll}}ation of a \uline{test functor}  A $\hand{\longrightarrow}$ (Cat)
(in contrast to a \uline{weak} test functor)  to functors satisfying the equiva-%
lent conditions (ii) and (iii),  implying that A is a test category (not
merely a weak one ! ),  and which seem stronger a priori than assuming
\deleted{moreover}\inserted{(besides A being a \deleted{weak} test category),} that i is a weak test functor  (condition (i)).  I confess I did
not make up my mind if \xout{this} for a test category A , there may be weak
test functors,  with all i(a) having final objects,  which are not test
functors in the present sense,  namely whether it may be true that \xout{all} for
\xout{categoriesx}\inserted[typed]{any C in (Cat)}  with final object,  $i^*(C)$ is aspheric,  and \xout{\illegible{1}} there-%
fore $i^*(\hand{\Delta}_1)$ aspheric,  without the \over{m}{l}atter being aspheric over $e_{\hand{A\hat{}}}$ i.e.
all products $i^*(\hand{\Delta}_1) \times a$  ($a \hand{\in}$ Ob A) being aspheric\inserted{(which would be automatic though \inserted{if A is even a \uline{strict} test category})}.  The difficulty here
is that it looks hard to check the condition for arbitrary C with final
object,  except through the homotopy argument using relative asphericity
of $i^*(\hand{\Delta}_1)$ .  Of course the distinction between weak test functors and test
functors is meaningful only as long as it is expected that the two are
actually distinct,  when applied to test categories.  Besides this,  the
restriction to categories i(a) with final objects looks theoretically a
little awkward,  and it should'nt be too hard I believe to get rid of it,
if need be.  But for the time being the only would-be test functors or
